\documentclass[a4paper,11pt]{ltxdoc}
\usepackage[UTF8,scheme=plain]{ctex}
\usepackage{geometry}
\usepackage{xcolor}
\usepackage{booktabs}
\usepackage{longtable}
\usepackage{array}
\usepackage{enumitem}
\usepackage{listings}
\usepackage{hypdoc}

\geometry{margin=2.55cm}
\setlength{\emergencystretch}{3em}
\hypersetup{
  colorlinks=true,
  linkcolor=blue!55!black,
  urlcolor=blue!55!black,
  citecolor=blue!55!black,
  pdftitle={Maki Notes Implementation Manual},
  pdfauthor={Maki Notes}
}

\definecolor{MakiBlue}{RGB}{42,91,142}
\definecolor{MakiGray}{RGB}{246,247,249}
\definecolor{MakiRule}{RGB}{210,216,224}

\lstdefinestyle{maki-doc}{
  basicstyle=\ttfamily\small,
  columns=fullflexible,
  keepspaces=true,
  frame=single,
  framerule=0.3pt,
  rulecolor=\color{MakiRule},
  backgroundcolor=\color{MakiGray},
  breaklines=true,
  tabsize=2
}
\lstset{style=maki-doc}

\newcommand{\mfile}[1]{\texttt{#1}}
\newcommand{\mcls}[1]{\textsf{#1}}
\newcommand{\mpkg}[1]{\textsf{#1}}
\newcommand{\mopt}[1]{\texttt{#1}}
\newcommand{\menv}[1]{\texttt{#1}}
\newcommand{\mcmd}[1]{\texttt{\textbackslash#1}}
\newcommand{\mkey}[1]{\texttt{#1}}
\newcommand{\Maki}{\textsf{Maki Notes}}

\newenvironment{implementation}[1]
  {\par\medskip\noindent\textbf{\color{MakiBlue}实现要点：#1}\par\smallskip}
  {\par\medskip}

\newenvironment{interfaceblock}[1]
  {\par\medskip\noindent\textbf{\color{MakiBlue}接口：#1}\par\smallskip\begin{description}[leftmargin=3.2cm,style=nextline]}
  {\end{description}\par\medskip}

\title{\Maki{} 实现手册\\
  \large 按功能划分的 \mfile{.cls}/\mfile{.sty} 源码说明}
\author{Maki Notes Project}
\date{2026-04-18}

\begin{document}
\maketitle
\tableofcontents

\begin{abstract}
本文是一份面向维护者的实现手册，而不是普通使用教程。它按照功能模块解释
\mfile{maki-notes.cls}、\mfile{maki-beamer.cls}、\mfile{maki-notes.sty}
和 \mfile{beamerthememaki.sty} 的职责边界、关键宏、选项流转、计数器策略、
主题/字体解析、讲义环境、研究笔记层、答案模式、TikZ 语义样式和构建流程。
写法参考 CTAN 包文档的组织方式：先给出文件映射与公共接口，再解释实现细节和扩展规则。
\end{abstract}

\section{文件架构}

\Maki{} 的源码被拆成两个文档类入口和两个样式实现层。

\begin{longtable}{@{}p{0.26\linewidth}p{0.68\linewidth}@{}}
\toprule
文件 & 职责 \\
\midrule
\mfile{maki-notes.cls} &
书式讲义入口。负责解析 \mopt{fontpreset}、\mopt{theme}、\mopt{answers}、
\mopt{solutions} 等类选项，装载 \mcls{ctexbook}，建立封面宏，再载入
\mfile{maki-notes.sty}。\\
\mfile{maki-beamer.cls} &
演示文稿入口。负责解析 Beamer 专属的 \mopt{layout}、\mopt{plain}、
\mopt{invert}、\mopt{accessibility}，装载 \mcls{ctexbeamer}，再把主题选项
转交给 \mfile{beamerthememaki.sty}。\\
\mfile{maki-notes.sty} &
共享样式主体。集中管理字体预设、主题色、数学排版、课程/研究元信息、
章节导读、研究条目导引、盒子环境、答案模式、引用命令、图文绕排、
页边批注和 TikZ/PGFPlots 语义样式。\\
\mfile{beamerthememaki.sty} &
Beamer 主题层。负责 title page、outline、headline、footline、进度条、
block 颜色和三种布局变体。\\
\bottomrule
\end{longtable}

\begin{implementation}{入口层和共享层分离}
两个 \mfile{.cls} 文件只做“选项解析 + 基类装配 + 少量顶层命令”。真正的视觉系统和内容环境放在
\mfile{maki-notes.sty}。这种结构保证讲义和 Beamer 可以共享同一套字体、主题色和 TikZ 样式。
\end{implementation}

\section{类选项和入口类}

\subsection{\mfile{maki-notes.cls}}

\mfile{maki-notes.cls} 使用 \mpkg{kvoptions} 声明类选项：

\begin{lstlisting}[language=TeX]
\DeclareStringOption[common]{fontpreset}
\DeclareStringOption[default]{theme}
\DeclareBoolOption[false]{answers}
\DeclareStringOption[inline]{solutions}
\DeclareDefaultOption{\PassOptionsToClass{\CurrentOption}{ctexbook}}
\ProcessKeyvalOptions*
\end{lstlisting}

\begin{interfaceblock}{讲义类选项}
\item[\mopt{fontpreset}] 取值为 \mopt{common}、\mopt{auto}、\mopt{windows}、\mopt{macos}、\mopt{linux}。
\item[\mopt{theme}] 取值为 \mopt{default}、\mopt{ocean}、\mopt{forest}、\mopt{graphite}、\mopt{amber}、\mopt{berry}、\mopt{sandstone}、\mopt{cobalt}、\mopt{sage}、\mopt{ruby}、\mopt{midnight}。
\item[\mopt{answers}] 布尔选项，控制是否显示答案。
\item[\mopt{solutions}] 字符串选项，控制答案输出位置：\mopt{inline}、\mopt{appendix}、\mopt{hidden}。
\end{interfaceblock}

解析完成后，类文件把选项值复制到共享层变量：

\begin{lstlisting}[language=TeX]
\edef\MakiFontPreset{\maki@fontpreset}
\edef\MakiTheme{\maki@theme}
\ifmaki@answers
  \MakiShowAnswerstrue
\else
  \MakiShowAnswersfalse
\fi
\edef\MakiSolutionsMode{\maki@solutions}
\end{lstlisting}

\subsection{\mfile{maki-beamer.cls}}

\mfile{maki-beamer.cls} 的结构与 notes 类相似，但基类换成 \mcls{ctexbeamer}，并增加 Beamer 主题选项。
类文件先把选项收集到 \mcmd{MakiBeamerThemeOptions}，再调用：

\begin{lstlisting}[language=TeX]
\expandafter\usetheme\expandafter[\MakiBeamerThemeOptions]{maki}
\end{lstlisting}

\begin{implementation}{为什么不用直接写死 \mcmd{usetheme}}
主题选项由类选项动态拼接，因此需要通过 \mcmd{expandafter} 把逗号分隔的选项列表展开后传给
\mcmd{usetheme}。这样用户只需要在类声明中写：
\begin{lstlisting}[language=TeX]
\documentclass[layout=splitbar,theme=cobalt]{maki-beamer}
\end{lstlisting}
不用再手动调用主题。
\end{implementation}

\section{共享样式层的模式检测}

\mfile{maki-notes.sty} 第一件事是判断自己运行在普通讲义还是 Beamer 中：

\begin{lstlisting}[language=TeX]
\newif\ifMakiBeamerMode
\MakiBeamerModefalse
\@ifclassloaded{beamer}{\MakiBeamerModetrue}{}
\@ifclassloaded{ctexbeamer}{\MakiBeamerModetrue}{}
\end{lstlisting}

随后很多包和环境都使用 \mcmd{ifMakiBeamerMode} 分支。例如：

\begin{itemize}
\item 讲义模式加载 \mpkg{geometry}、\mpkg{fancyhdr}、\mpkg{caption}、\mpkg{wrapfig}、\mpkg{marginnote}。
\item Beamer 模式跳过页边距、页眉页脚、浮动题注等书式文档功能。
\item 研究环境和 TikZ 样式保留给两种模式共用。
\end{itemize}

\begin{implementation}{模式分支的目的}
共享层不是一个完全通用的 LaTeX 包，而是给 \mcls{maki-notes} 和 \mcls{maki-beamer} 使用的样式系统。
模式检测避免了 \mcls{beamer} 与 \mpkg{geometry}/\mpkg{fancyhdr}/\mpkg{caption} 等包之间的冲突。
\end{implementation}

\section{字体预设系统}

字体系统分成拉丁字体和 CJK 字体两层。

\subsection{拉丁字体}

\mcmd{makisetlatinfonts} 优先使用 TeX Gyre Pagella、TeX Gyre Heros 和 Latin Modern Mono。
这些字体通常随 TeX Live 安装，能保证跨平台可用性。

\subsection{CJK 字体}

CJK 预设由一组 \mcmd{makiapplyfontpreset...} 宏实现：

\begin{longtable}{@{}p{0.36\linewidth}p{0.56\linewidth}@{}}
\toprule
宏 & 实现策略 \\
\midrule
\mcmd{makiapplyfontpresetcommon} & 使用 Fandol 宋体、黑体、仿宋、楷体，作为 TeX Live 默认安全方案。\\
\mcmd{makiapplyfontpresetwindows} & 检测 \mopt{SimSun}，再选择微软雅黑、黑体、楷体等 Windows 常见字体。\\
\mcmd{makiapplyfontpresetmacos} & 检测 \mopt{Songti SC}，再选择苹方、黑体、楷体等 macOS 字体。\\
\mcmd{makiapplyfontpresetlinux} & 检测 Noto/Source Han CJK 字体，不存在时回退到 \mopt{common}。\\
\mcmd{makiapplyfontpresetauto} & 依次探测 macOS、Windows、Linux 字体，最后回退到 \mopt{common}。\\
\bottomrule
\end{longtable}

\begin{implementation}{字体选择的容错}
\mcmd{makipickfirstfont} 接收一个逗号分隔候选列表，使用 \mcmd{IfFontExistsTF} 找到第一个可用字体。
如果没有命中，则写入回退字体。显式请求的平台预设不存在时，样式层使用
\mcmd{PackageWarningNoLine} 给出警告并退回 \mopt{common}。
\end{implementation}

\section{主题色系统}

主题通过 \mcmd{makiapplytheme...} 宏实现。每个主题都会定义相同的一组语义色位：

\begin{itemize}
\item \mkey{LogicColor}/\mkey{LogicBg}：定理、证明、强调逻辑。
\item \mkey{DefColor}/\mkey{DefBg}：定义、公理、输入输出。
\item \mkey{ExColor}/\mkey{ExBg}：例题、练习、流程主线。
\item \mkey{NoteColor}：注释、方法、研究想法。
\item \mkey{GeometryColor}/\mkey{GeometryBg} 与 \mkey{TopologyColor}/\mkey{TopologyBg}：几何、拓扑、图形语义。
\item \mkey{NNInputColor}、\mkey{NNHiddenColor}、\mkey{NNOutputColor}、\mkey{NNTensorColor}：神经网络/Tensor 图形语义。
\item \mkey{MainText}：正文主色。
\end{itemize}

\begin{implementation}{主题切换不改公开接口}
盒子环境、TikZ 样式和 Beamer 主题都引用语义色位，而不是硬编码某个主题的 RGB 值。
因此新增主题只需要添加一个 \mcmd{makiapplytheme<name>} 宏，并在文档中登记合法取值。
\end{implementation}

\section{数学排版微调}

\mfile{maki-notes.sty} 对一些数学符号做了讲义化调整：

\begin{itemize}
\item \mcmd{ge}、\mcmd{le} 改为 \mcmd{geqslant}、\mcmd{leqslant}。
\item \mcmd{lim}、\mcmd{int}、\mcmd{sum} 默认使用 display style，适配中文讲义的大字号阅读。
\item \mcmd{dfrac} 增加不可见高度/深度 strut，减少分式上下拥挤。
\item \menv{pmatrix}、\menv{bmatrix}、\menv{dcases} 等环境通过 LaTeX hook 调整行距。
\end{itemize}

\begin{implementation}{风险边界}
这些改写会改变全局数学外观，因此它们适合本模板的讲义定位。如果要把 \mfile{maki-notes.sty}
拆成通用 CTAN 包，应考虑提供关闭这些重定义的选项。
\end{implementation}

\section{课程元信息和章节导读}

\subsection{键值接口}

课程、章节和导读信息使用 \mpkg{pgfkeys} 解析：

\begin{lstlisting}[language=TeX]
\SetCourseInfo{
  course={高等数学},
  term={2026 春},
  lecture={第 3 讲},
  date={2026-04-13}
}
\end{lstlisting}

内部键族包括：

\begin{longtable}{@{}p{0.3\linewidth}p{0.62\linewidth}@{}}
\toprule
键族 & 存储变量 \\
\midrule
\mkey{/maki/course} & \mcmd{MakiCourseInfoCourse}、\mcmd{MakiCourseInfoTerm}、\mcmd{MakiCourseInfoLecture} 等。\\
\mkey{/maki/chapter} & 暂存到 \mcmd{MakiPendingChapter...}，在下一次 chapter 时绑定。\\
\mkey{/maki/guide} & 存储本章路线、核心公式、方法、误区、考试关注等导读内容。\\
\bottomrule
\end{longtable}

\subsection{章节绑定}

讲义模式中使用 \mcmd{pretocmd} 挂钩 \mcmd{@chapter} 和 \mcmd{@schapter}：

\begin{lstlisting}[language=TeX]
\pretocmd{\@chapter}{\MakiBindPendingChapterInfo}{}{}
\pretocmd{\@schapter}{\MakiBindPendingChapterInfo}{}{}
\end{lstlisting}

\mcmd{MakiBindPendingChapterInfo} 会把“下一章”的待绑定信息复制到当前章变量，并清空 pending 状态。
这样用户可以在 \mcmd{chapter} 前写 \mcmd{SetChapterInfo}，不会污染后续章节。

\subsection{导读页生成}

\mcmd{MakeChapterGuide} 负责生成整页导读：

\begin{itemize}
\item 章节标题来自 \mcmd{chaptermark} 写入的 \mcmd{MakiCurrentChapterTitle}。
\item 课程信息来自 \mcmd{SetCourseInfo}。
\item 学习目标、路线、公式、误区等字段通过 \menv{tcbraster} 排成卡片。
\item 局部目录通过 \mpkg{etoc} 的 \mcmd{localtableofcontents} 生成。
\end{itemize}

\section{研究笔记工作流}

研究层与课程层并行，使用 \mcmd{SetResearchInfo}、\mcmd{SetEntryInfo} 和
\mcmd{MakeEntryGuide} 维护项目上下文、条目上下文和条目导引块。

\begin{interfaceblock}{研究接口}
\item[\mcmd{SetResearchInfo}] 设置项目、方向、作者、合作者、状态、创建时间、更新时间和标签。
\item[\mcmd{SetEntryInfo}] 设置当前条目的主题、来源、source key、更新时间、标签和目标。
\item[\mcmd{MakeEntryGuide}] 输出条目导引卡片和摘要/依赖/目标/障碍/下一步等字段。
\item[\mcmd{DeclareSymbol}] 登记符号说明。
\item[\mcmd{PrintSymbolIndex}] 输出符号表。
\item[\mcmd{PrintOpenProblemIndex}] 输出开放问题表。
\end{interfaceblock}

\begin{implementation}{动态字段存储}
符号表和开放问题表使用计数器 + \mcmd{csname} 动态控制序列存储条目。例如第 \(n\) 个符号会存入
\verb|\maki@symbol@n@symbol| 和 \verb|\maki@symbol@n@desc|。打印时用 \mcmd{@whilenum}
从 1 遍历到当前计数器值。
\end{implementation}

\section{图文绕排和页边批注}

讲义模式额外加载 \mpkg{wrapfig}、\mpkg{marginnote}、\mpkg{ifoddpage}，提供：

\begin{interfaceblock}{图文/页边接口}
\item[\menv{leftfiguretext}] 左图右文。
\item[\menv{rightfiguretext}] 右图左文。
\item[\menv{figuretextleft}] \menv{leftfiguretext} 的别名式封装。
\item[\menv{figuretextright}] \menv{rightfiguretext} 的别名式封装。
\item[\mcmd{sidenote}] 无标题页边批注。
\item[\mcmd{marginremark}] 带可选标题的页边批注。
\item[\mcmd{sidecaption}] 使用页边批注输出 figure 题注。
\end{interfaceblock}

\begin{implementation}{双面排版}
\mcmd{maki@outernote} 通过 \mcmd{checkoddpage} 判断奇偶页。在双面模式下，奇数页使用右侧格式，
偶数页使用左侧格式；单面模式固定在右侧。
\end{implementation}

\section{TikZ 和 PGFPlots 语义样式}

模板把 TikZ 样式按语义分组，而不是只提供颜色片段。

\begin{longtable}{@{}p{0.28\linewidth}p{0.64\linewidth}@{}}
\toprule
样式族 & 用途 \\
\midrule
\mkey{nn ...} & 神经网络、张量、卷积块、跳连、注意力连接。\\
\mkey{transformer ...} & token、attention block、FFN、norm、residual flow 等 Transformer 结构图。\\
\mkey{pinn ...}/\mkey{fpinn ...} & PINN/fPINN 的物理约束、边界、初值、loss、分数阶算子等。\\
\mkey{math ...}/\mkey{geom ...} & 坐标轴、曲线、切线、法线、几何点和角标。\\
\mkey{topo ...}/\mkey{manifold ...}/\mkey{sphere ...} & 拓扑基本域、流形局部图、三维球面。\\
\mkey{fourier ...} & 时域、频域、窗函数、变换箭头和谱线。\\
\mkey{flow ...}/\mkey{timeline ...} & 流程图和时间轴。\\
\mkey{prob ...}/\mkey{graph ...} & 概率集合图和图论节点边图。\\
\bottomrule
\end{longtable}

\mcmd{MakiFitTikZ} 使用 savebox 测量内容宽度，只在超过目标宽度时调用 \mcmd{resizebox}：

\begin{lstlisting}[language=TeX]
\sbox{\MakiFitTikZBox}{#2}
\ifdim\wd\MakiFitTikZBox>#1\relax
  \resizebox{#1}{!}{\usebox{\MakiFitTikZBox}}
\else
  \usebox{\MakiFitTikZBox}
\fi
\end{lstlisting}

这比无条件缩放更稳：较小的图保持原始线宽和文字大小，只有过宽图才收缩。

\section{盒子环境和计数器}

讲义环境基于 \mpkg{tcolorbox}。基础样式分三层：

\begin{enumerate}
\item \mkey{commonbox}：定理、定义、例题等核心盒子的通用外观。
\item \mkey{maki workflow card}/\mkey{maki workflow hero}：章节导读和研究导引卡片。
\item \mkey{maki workflow block}：核心公式、方法、误区、研究对象、答案等信息块。
\end{enumerate}

\begin{interfaceblock}{主要环境}
\item[\menv{theorem}, \menv{lemma}, \menv{proposition}, \menv{corollary}] 逻辑类盒子。
\item[\menv{definition}, \menv{axiom}] 定义类盒子。
\item[\menv{example}, \menv{exercise}, \menv{realexam}] 例题、练习、真题。
\item[\menv{notation}, \menv{question}, \menv{conjecture}, \menv{claim}] 研究对象。
\item[\menv{keyformula}, \menv{pitfall}, \menv{methodnote}, \menv{examfocus}] 教学工作流信息块。
\item[\menv{remark}, \menv{proofsketch}, \menv{idea}, \menv{obstacle}, \menv{nextstep}] 研究过程记录块。
\end{interfaceblock}

\begin{implementation}{共享 theorem 编号}
\menv{theorem}、\menv{definition}、\menv{example}、\menv{exercise} 等讲义盒子在当前实现中通过
下面的 \mpkg{tcolorbox} 选项推进同一个 \mkey{theorem} 计数器：
\begin{lstlisting}[language=TeX]
before title={\refstepcounter{theorem}}
\end{lstlisting}
这会形成统一的“讲义对象编号流”。源码中也声明了 \mkey{definition} 和 \mkey{example}
计数器，但当前盒子标题没有使用它们输出编号。
\end{implementation}

\section{答案模式}

答案模式由 \mfile{maki-notes.cls} 解析选项，再由 \mfile{maki-notes.sty} 实现。

\begin{interfaceblock}{答案接口}
\item[\mopt{answers=false}] 默认值，不输出 \menv{solution} 内容。
\item[\mopt{answers=true,solutions=inline}] 在原位置输出答案盒子。
\item[\mopt{answers=true,solutions=appendix}] 登记答案，等 \mcmd{PrintSolutions} 统一输出。
\item[\mopt{solutions=hidden}] 即使打开 \mopt{answers=true} 也隐藏答案。
\end{interfaceblock}

\begin{implementation}{收集式答案}
\menv{solution} 使用 \mcmd{NewDocumentEnvironment} 的 \verb|+b| 参数捕获环境正文。
附录模式下，正文不立即排版，而是存入
\verb|\maki@solution@<n>@title| 和 \verb|\maki@solution@<n>@body|。调用
\mcmd{PrintSolutions} 时再遍历输出。
\end{implementation}

\begin{implementation}{Beamer 兼容}
\mcls{beamer} 自身已经定义过 \menv{solution} theorem。为了避免冲突，本模板的答案收集系统只在
非 Beamer 模式启用。
\end{implementation}

\section{引用辅助命令}

模板提供一组轻量中文引用包装：

\begin{lstlisting}[language=TeX]
\chapref{chap:intro}
\secref{sec:setup}
\thmref{thm:main}
\figref{fig:diagram}
\eqnref{eq:energy}
\end{lstlisting}

这些命令只是 \mcmd{ref}/\mcmd{eqref} 的稳定包装，不引入 \mpkg{cleveref}。好处是包依赖少、
行为可预测；代价是不会自动识别 label 类型，用户需要选择对应命令。

\section{Beamer 主题实现}

\mfile{beamerthememaki.sty} 通过 Beamer 主题接口提供三套布局：

\begin{longtable}{@{}p{0.22\linewidth}p{0.7\linewidth}@{}}
\toprule
布局 & 实现内容 \\
\midrule
\mopt{durham} & 默认布局。标题页右下角有主题色三角块，headline 显示横向章节导航，footline 显示页码与进度条。\\
\mopt{minimal} & 减少装饰，使用细竖条和简化页眉页脚，适合内容密集型 slides。\\
\mopt{splitbar} & 顶部分区明显，outline 页使用左侧色块，适合课程式结构。\\
\bottomrule
\end{longtable}

布局选择使用 \mcmd{csname} 调度：

\begin{lstlisting}[language=TeX]
\setbeamertemplate{title page}{\csname MakiBeamerTitlePage@\MakiBeamerLayout\endcsname}
\setbeamertemplate{headline}{\csname MakiBeamerHeadline@\MakiBeamerLayout\endcsname}
\setbeamertemplate{footline}{\csname MakiBeamerFootline@\MakiBeamerLayout\endcsname}
\end{lstlisting}

\begin{implementation}{主题色复用}
Beamer 主题不重新定义完整配色，而是读取共享层定义的 \mkey{LogicColor}、\mkey{ExColor}、
\mkey{DefColor}、\mkey{MainText} 等语义色位。若用户单独加载主题，主题文件会提供最低限度的 fallback 色。
\end{implementation}

\section{构建、测试和发布}

\mfile{Makefile} 维护常用构建入口：

\begin{longtable}{@{}p{0.22\linewidth}p{0.7\linewidth}@{}}
\toprule
目标 & 说明 \\
\midrule
\mkey{make main} & 编译 \mfile{document2.tex}。\\
\mkey{make example} & 编译 \mfile{example.tex}。\\
\mkey{make beamer-demo} & 编译 \mfile{beamer-demo.tex}。\\
\mkey{make tikz-templates} & 编译 \mfile{tikz-template-pages.tex}。\\
\mkey{make smoke} & 编译最小打包测试文档。\\
\mkey{make test} & 编译 smoke 和所有测试文档。\\
\mkey{make all} & 编译主文档、示例、Beamer、TikZ 模板册和测试。\\
\bottomrule
\end{longtable}

\noindent\mfile{.github/workflows/test.yml}
在 push/PR 时编译 smoke 和测试文档。

\noindent\mfile{.github/workflows/release.yml}
在 tag push 时编译完整示例、打包源码并上传 release assets。

\section{扩展指南}

\subsection{新增主题}

\begin{enumerate}
\item 在 \mfile{maki-notes.sty} 中新增 \mcmd{makiapplytheme<name>}。
\item 定义完整语义色位，不要只覆盖部分颜色。
\item 在 README、template guide 和测试文档中登记取值。
\item 用 \mfile{tests/test-tikz-styles.tex} 验证 TikZ 样式跟随主题变化。
\end{enumerate}

\subsection{新增讲义环境}

\begin{enumerate}
\item 先决定是否应使用共享 \mkey{theorem} 计数器。
\item 在 \mcmd{tcbset} 中添加语义盒子样式，或复用现有 \mkey{maki workflow block}。
\item 用 \mcmd{NewDocumentEnvironment} 或 \mcmd{newenvironment} 暴露用户接口。
\item 新增测试文档或扩展现有测试文档。
\end{enumerate}

\subsection{新增 TikZ 样式族}

\begin{enumerate}
\item 按语义命名，例如 \mkey{flow process}、\mkey{graph edge strong}。
\item 颜色引用语义色位，不直接绑定某个主题。
\item 在 \mfile{tikz-template-pages.tex} 中加入可复制模板页。
\item 在 \mfile{tests/test-tikz-styles.tex} 中加入最小覆盖图。
\end{enumerate}

\section{维护注意事项}

\begin{itemize}
\item 不要让 \mfile{.cls} 直接实现大量视觉逻辑；入口类应保持薄层。
\item 新增公开选项时，先在类文件解析，再传给 \mfile{maki-notes.sty} 或主题层处理。
\item 涉及 Beamer 的改动必须确认 \mcmd{ifMakiBeamerMode} 分支。
\item 涉及浮动、caption、wrapfig、fancyhdr 的改动通常只适合 notes 模式。
\item 新增计数器或索引时，应同步考虑 \mcmd{label}/\mcmd{ref} 行为和测试覆盖。
\item 完成修改后至少运行 \verb|make -B test|；发布前运行 \verb|make -B all|。
\end{itemize}

\end{document}
